% Proposed isolated A.88. All checkpoint28 sources remain unchanged.
\statement{A.88}{Proved}{A full-head Alexander coefficient test and the canonical congruence modulo eight}
Retain the actual negative-GT0 data and the two full units
\(a=a_{\rm reg}\), \(d=d_\alpha=2k+1\) of A.86--A.87.
Let \(K=\langle\!\langle x\rangle\!\rangle_F\), let
\(A=F/K=\widehat{\mathbb Z}\) with generator T the image of y, and put
\[
 \begin{gathered}
 \Lambda=\mathbb F_2[[A]],\qquad
 K/\overline{K^2[K,K]}=\Lambda[x],\qquad U=[u]\in\Lambda,\\
 \overline P(T)=P(T^{-1}),\qquad
 \Theta_u=1+(T-1)(U+\overline U).
 \end{gathered}
 \tag{A.88.1}
\]
The bar on a coefficient denotes the displayed inversion, not a
subgroup closure. For \(b\in\widehat{\mathbb Z}\), let \(D_b\in\Lambda\)
be the unique continuous geometric sum defined by
\[
                         (T-1)D_b=T^b-1.
 \tag{A.88.2}
\]
Then the actual full-head identity is
\[
             \Theta_uD_{a_{\rm reg}}=1.
 \tag{A.88.3}
\]
In particular the remaining full coefficient equality has the exact
test using the canonical u alone:
\[
              a_{\rm reg}=d_\alpha
              \quad\Longleftrightarrow\quad\Theta_uD_{d_\alpha}=1
              \quad\hbox{in }\mathbb F_2[[\widehat{\mathbb Z}]].
 \tag{A.88.4}
\]
The head in this ring is the full procyclic group, so (A.88.4) is
not only a dyadic comparison. Moreover the actual units satisfy
\[
                         a_{\rm reg}\equiv d_\alpha\pmod 8.
 \tag{A.88.5}
\]
Equivalently A.87's two canonical frame coordinates satisfy
\(t\equiv\rho_D\pmod2\). Consequently
\[
           (\Omega_\alpha)_{\rm ab}
              \in4\widehat{\mathbb Z}\,(1,-1).
 \tag{A.88.6}
\]
No higher congruence, full equality of the two units, full Fox-unit
claim, or nonlinear vanishing of \(\Omega_\alpha\) is inferred.

\begin{proof}
\emph{1. The actual full module and its semilinear actions.}\par
The Alexander identification in (A.88.1) is C.44's exact cyclic-cover
calculation with the two basis labels exchanged. To retain its full
scope explicitly, the subgroup
\(U_n=\ker(F\to\mathbb Z/n)\), with x sent to zero and y to one,
has Schreier basis
\[
                   y^n,\qquad y^i x y^{-i}\quad(0\leq i<n).
\]
For \(n\mid n'\) the inclusion on mod-two homology multiplies the
first coordinate by \(n'/n\) and reduces the boundary indices
modulo n. Passage from n to2n kills the first coordinate in the
inverse limit, leaving exactly \(\mathbb F_2[[\widehat{\mathbb Z}]]\).
Every continuous character of K extends to an open neighborhood KV:
choose an open normal V with \(K\cap V\) in its kernel and assign
zero on V. Commutators \([V,K]\) lie in that kernel, so this is
well-defined. Every open neighborhood of K contains some \(U_n\).
Dualizing the resulting direct limit of character groups proves
that the inverse limit is exactly the stated quotient of K.
Conjugation by y shifts the indices, giving the module action.

Multiplication by \(T-1\) is an isomorphism of compact modules
from \(\Lambda\) to its augmentation ideal. Its compact image
surjects onto the augmentation ideal at every finite cyclic head.
An element of its kernel has constant coefficients at each such
head; the transition from the2n head to the n head sends the
constant vector to zero. Thus the kernel is zero. This proves
existence, uniqueness and continuous dependence in (A.88.2), as
well as the ordinary geometric-sum formula on nonnegative integers.

Both the actual alpha and \(j_\delta=\operatorname{Ad}(\delta^{-1})|_F\)
preserve K and invert A. On its module they send [x] respectively
to \(-[x]=[x]\) and to [x]. Therefore each acts by the precise
semilinearity \(P(T)\mapsto P(T^{-1})\). The extra conjugator f
in \(\alpha(y)=f^{-1}y^{-1}f\) lies in K and has trivial inner
action on this abelian module. Since \(u,f\in F'\subset K\), the
actual canonical equation \(u=f\alpha(u)\) gives
\[
                          [f]=U+\overline U.
 \tag{A.88.7}
\]
Its right side is invariant under the coefficient bar. The
augmentation of U is the x-exponent of u modulo two, hence zero.
There is a unique \(V\in\Lambda\) such that
\[
                              U=(T-1)V.
 \tag{A.88.8}
\]

\emph{2. Ordered derivation of the full identity.}\par
A.87 gives the existing normalized regular action
\[
 \sigma_x=j_\delta\alpha,\qquad
 \sigma_x(x)=x^{-1},\qquad
 \sigma_x(y)=j_\delta(f)^{-1}xyj_\delta(f).
\]
It is the identity on A and on the mod-two module, since the two
head inversions cancel. Set \(\kappa=\sigma_x(y)y^{-1}\). Passing
the displayed word to the module, with its factor order retained,
gives
\[
 \begin{aligned}
 [\kappa]
   &=[j_\delta(f)]+1+T[j_\delta(f)]\\
   &=1+(T-1)(U+\overline U)=\Theta_u.
 \end{aligned}
 \tag{A.88.9}
\]
The coefficient T comes from conjugating the final factor by y;
the two signs coincide only after passage to characteristic two.

The unique actual solution G of A.87 satisfies
\(\sigma_x(G)G^{-1}=x\) and has image \(y^a\) in A. Write
\(G=k_Gy^a\) with \(k_G\in K\). The class contributed by \(k_G\)
cancels because \(\sigma_x\) acts as the identity on the module.
For integer powers the crossed-homomorphism formula gives
\([\sigma_x(y^b)y^{-b}]=\Theta_uD_b\); continuity proves the same
for every profinite b. The equation for G now gives (A.88.3).
In particular \(\Theta_u\) is a unit in this commutative ring.

If \(\Theta_uD_d=1\), multiplication by its inverse gives
\(D_d=D_a\); (A.88.2) then gives \(T^d=T^a\). The full completed
group algebra distinguishes these group elements: any distinct
exponents have distinct residues in some finite cyclic head, where
the corresponding group elements are different basis vectors over
\(\mathbb F_2\). Hence d=a. The reverse implication is (A.88.3).
This proves (A.88.4) without discarding any prime in the group head.

\emph{3. The first nontrivial augmentation coefficient.}\par
Project the full ring to
\[
                \mathbb F_2[C_4]=\mathbb F_2[\varepsilon]/(\varepsilon^4),
                \qquad\varepsilon=T-1.
\]
Write the image of V as
\(v_0+v_1\varepsilon+v_2\varepsilon^2+v_3\varepsilon^3\).
Then \(v_0\) is exactly the augmentation of V and
\[
 \begin{gathered}
 U=v_0\varepsilon+v_1\varepsilon^2+v_2\varepsilon^3,\\
 \overline\varepsilon=\varepsilon+\varepsilon^2+\varepsilon^3,
 \qquad (\overline{\varepsilon})^2=\varepsilon^2,
 \qquad (\overline{\varepsilon})^3=\varepsilon^3.
 \end{gathered}
\]
Thus
\[
                 U+\overline U=v_0\varepsilon^2+v_0\varepsilon^3,
                 \qquad\Theta_u=1+v_0\varepsilon^3.
 \tag{A.88.10}
\]
The actual \(a=1-4t\) is1 modulo4. In this finite ring the
geometric sums obey
\(D_{b+c}=D_b+T^bD_c\), with
\(D_4=\varepsilon^3\) and \(D_8=0\). Consequently for
\(a=1+4c\in\widehat{\mathbb Z}\), continuity from the integers gives
\[
                     D_a=1+(c\bmod2)\varepsilon^3.
 \tag{A.88.11}
\]
This calculation uses the continuous D-function in the full ring;
it does not divide by the zero divisor \(\varepsilon\) in the
finite quotient. Substitution in (A.88.3) gives
\[
                         (a-1)/4\equiv\operatorname{aug}(V)\pmod2.
 \tag{A.88.12}
\]

\emph{4. Compare through the same actual dihedral marking.}\par
Set
\[
 R=\mathbb F_2[\varepsilon]/(\varepsilon^2),\qquad
 \mathcal D_8=R_{\rm add}\rtimes C_2,
\]
where the nontrivial element of \(C_2\) multiplies R by
\(T=1+\varepsilon\). Send the original free generators to
\[
                          x\longmapsto(1,0),\qquad
                          y\longmapsto(0,1).
 \tag{A.88.13}
\]
Both images are involutions. Their product has square
\((\varepsilon,0)\ne1\), and they generate the whole group, since
1 and its T-translate span R. Thus (A.88.13) is precisely an
order-eight dihedral quotient of the original x,y marking, with
rotation \(T_D=xy\) of order four.

On K, (A.88.13) is reduction of the exact Alexander coefficient
to R: conjugation by y multiplies the image of x by T. This checks
the assertion on every normal generating conjugate of x and hence
continuously on K. The canonical u therefore has image
\[
                  (U\bmod\varepsilon^2,0)
                    =(\operatorname{aug}(V)\varepsilon,0).
\]
Its A.45 dihedral image, on the other hand, is \(T_D^k\), with k
even. At the same order-eight quotient this is
\(((k/2\bmod2)\varepsilon,0)\). We conclude that
\[
                         \operatorname{aug}(V)\equiv k/2\pmod2.
 \tag{A.88.14}
\]
This is a common quotient of the same full canonical data, not
an independently chosen ring example or a finite target census.

Combining (A.88.12)--(A.88.14) with \(d=2k+1\) yields
\((a-1)/4\equiv(d-1)/4\pmod2\), proving (A.88.5).
The divisions are defined by the previously proved divisibility
and are unique because \(\widehat{\mathbb Z}\) is torsion-free.
Now A.86 gives
\((\Omega_\alpha)_{\rm ab}=((d-a)/2,-(d-a)/2)\), which proves
(A.88.6). The relation \(t\equiv\rho_D\pmod2\) follows from
\(a=1-4t\) and \(d=1-4\rho_D\).

The full coefficient gate is still (A.88.4); (A.88.5) is only
its stated dyadic consequence. In particular the original marked
R2 input and the nonlinear equation of A.87 remain separate
requirements. No unit equality or full marked conjugator is
deduced from a comparison at this one common finite quotient.
\end{proof}
